\color{unidarkgreen}{\subsection{Газовая защита}}
\color{black}

\begin{enumerate}[label=\arabic{section}.\arabic{subsection}.\arabic*,  align=left, labelsep=4pt, leftmargin=0pt, itemindent=57pt]

\item
Газовая защита предназначена для защиты маслонаполненного оборудования от всех видов
внутренних повреждений, сопровождающихся выделением газа, интенсивным протеканием масла из
бака в расширитель, а также от утечки масла.

\item
Логикой предусматривается прием сигналов от:
\begin{itemize}
\item сигнального газового реле бака ШР;
\item отключающего газового реле бака ШР.
\end{itemize}

\item
Для повышения надежности работы газовой защиты предусмотрена возможность резервирования
сигналов срабатывания газовых защит и сигналов от устройств контроля изоляции (тока утечки) цепей
газовых защит путем их приема от двух источников -- ПДС1 и ПДС2. Устройства контроля изоляции,
установленные в цепях газовых защит, при снижении изоляции в течение времени, превышающем
значение параметра \textbf{«Тбл ГЗ ШР откл.»} формируют сигнал о неисправности соответствующих цепей
газовой защиты, который в свою очередь может действовать на блокировку сигналов срабатывания
газовой защиты.

\item
Предусмотрена возможность перевода действия сигнальных ступеней газовой защиты на
отключение, которая может быть использована в режиме опробования трансформатора. В случае
обслуживания трансформатора (например, доливке масла) предусмотрен перевод действия
отключающих ступеней газовой защиты на сигнал.

\begin{figure}[H]
\centering
\includegraphics[width=0.7\textwidth,height=0.7\textheight,keepaspectratio]{\fbpath/161.1073 ГЗ/img/GZ_sign.pdf}
\captionof{figure}{Функциональный блок ГЗ ШР сигнальная ступень}\label{pict_gz:GZ_sign}
\end{figure}

\begin{figure}[H]
\centering
\includegraphics[width=0.7\textwidth,height=0.7\textheight,keepaspectratio]{\fbpath/161.1073 ГЗ/img/GZ_otkl.pdf}
\captionof{figure}{Функциональный блок ГЗ ШР отключающая ступень}\label{pict_gz:GZ_otkl}
\end{figure}

\small
\begin{longtable}{|>{\centering\arraybackslash}m{5.3cm}|>{\centering\arraybackslash}m{4.2cm}|>{\centering\arraybackslash}m{1cm}|>{\centering\arraybackslash}m{3.3cm}|>{\centering\arraybackslash}m{1.8cm}|}
% Первая страница: полный заголовок
\caption{Перечень уставочных параметров ГЗ\hfill\vspace{-0.5\baselineskip}}\label{gz:tbl1}\\
\hline
\rowcolor{gray!30}
Наименование уставки & Диапазон & Шаг & Значение по умолчанию & Единица измерения \\
\hline
\endfirsthead

% Промежуточные страницы: «Продолжение таблицы»
\caption*{\hspace{3pt}\emph{Продолжение таблицы \ref{gz:tbl1}}\hfill\ }\vspace{-0.5\baselineskip} \\
\hline
\rowcolor{gray!30}
Наименование уставки & Диапазон & Шаг & Значение по умолчанию & Единица измерения \\
\hline
\endhead

%===m>ГЗ|-
%===m
\end{longtable}
\normalsize

\end{enumerate}